\section{R4: Indexed families and motive search}\label{sec:r4}

\question{R4.Q1}{Can a useful motive class be enumerated completely and terminate?}
\status{Yes at a fixed template bound; index-occurrence abstraction alone is not closed under accumulator generalization.}

A practical motive search should start from the recursor's dependent telescope,
not from an untyped enumeration of expressions of type \code{Sort}. Compute
candidate abstractions of selected, maximal occurrences of indices and the
major premise in the goal. Abstract consistently, preserve dependencies, and
reject any candidate whose motive type is wrong. If there are $k$ eligible
occurrences, a raw subset enumeration is finite, with at most $2^k$ choices;
typed equivalence and repeated occurrences can reduce it substantially.

That finite class does not automatically include accumulating recursion. A
motive like $M(n,x)=S\to T(n,x)$ contains a new binder not created by replacing
an index occurrence. Close the template language explicitly under a bounded
number of dependent products, product-valued strengthening, and optional
initialization state. Enumerate dependency-closed subsets of existing locals
for generalization and bounded fresh auxiliary types from the carrier grammar.
This gives an exact, useful answer to the question: completeness relative to
the defined template language, not completeness for arbitrary motives.

For \code{Vec A n -> Vec B n}, the intended motive abstracts the index and
vector to yield \code{Vec B n}; no accumulator invention is necessary. Preserve
fixed parameters $A$, $B$, and the mapping function, and give the recursive
branch its correctly indexed induction hypothesis. Generalizing every local
would unnecessarily enlarge the branch interface. Conversely, an accumulator
that changes at recursive calls must be generalized even when its type does
not depend on an index. Syntactic index-dependency closure alone will miss it.

\question{R4.Q2}{How should transports be represented without destroying canonical enumeration?}
\status{Canonicalize index expressions where justified, retain typed equality evidence, and separate display erasure from kernel terms.}

Proof irrelevance concerns inhabitants of the same proposition. It does not
turn every pair of propositionally equal indices into definitionally equal
indices, nor does it make every \code{Eq.rec} computation disappear. A term of
$\mathsf{Vec}\ A\ n$ cannot simply be inserted where
$\mathsf{Vec}\ A\ m$ is expected because some proof of $n=m$ exists somewhere.
The equality and the appropriate transport must actually be supplied.

Use a small index-normalization service. It returns a normalized index and a
proof of equality to the original, or reports that no supported normalization
applies. Constructor arithmetic and selected associativity/identity rules can
be handled predictably. Arbitrary \code{simp} search can be an optional proof
service, but its success is not a definition of a globally canonical form.
The example $\mathsf{Vec}\ A\ (n+0)$ needs no propositional cast with Lean's
usual addition; $\mathsf{Vec}\ A\ (0+n)$ at symbolic $n$ is the instructive
case \cite{leanprelude}.

A typed congruence structure may store equivalent expressions together with
proofs and their types. Heterogeneous equality can help express relationships
before source and destination types coincide, but it does not remove the need
for a typed reconstruction theorem. Replacing Lean's typing problem by a
setoid problem would move rather than solve the difficulty.

Rank candidates using a proof-insensitive representation only after stating
which equalities it respects. Proof subterms may have zero semantic ranking
cost, while transport nodes still have construction and checking costs.
A displayed term with omitted casts must be elaborated again and proved equal
to the certified term, or be explicitly labeled as a presentation rather than
paste-ready Lean. Bound transports to control search, but record that bound;
it restricts the grammar and may exclude valid programs.

\question{R4.Q3}{What unfolding strategy is supported by existing experience?}
\status{Use staged, demand-driven unfolding in Leant2, informed by actual tool behavior rather than a universal transparency rule.}

Canonical's published encoding unfolds definitions that alias function types,
including \code{Set} and \code{Not}, and unfolds reducible definitions and
type-class instances before constructing its search problem. This is more
specific than describing it simply as fixed transparency or complete
normalization. It also supports generalized induction over indexed types,
so it is a stronger baseline than the proposal's comparison suggests
\cite{canonical}.

The sensible adaptation is a staged policy. First expose only the head shape
needed by the next rule: a function binder for introduction, an inductive head
for elimination, or a class head for instance resolution. Use a small,
predictable reducible view for indexing. Escalate transparency only when a
particular matching or conversion attempt is stuck, and charge the escalation
to that attempt. Cache under the transparency mode and relevant environment
fingerprint. Do not normalize every program and every local type fully at
every node.

The official \code{sauto} documentation gives a concrete comparison:
\code{unfold:} names definitions to unfold heuristically, with an empty list by
default; \code{unfold!:} forces the listed unfoldings. Logical \code{iff} and
\code{not} receive special treatment, while \code{red} simplification and its
\code{ered} eager scheduling are enabled by default. This is a configurable
mixture of reduction and selective unfolding, not blanket normalization of
all definitions \cite{sauto}. Pin the version and these options in any
comparison; the documented defaults are not a theorem about every historical
implementation path.

A useful ablation measures head exposure, reducible normalization, and
aggressive unfolding separately. Count conversion failures, expansions,
normalization time, and cache growth. A transparency policy that saves a few
unifications but expands huge definitions is not necessarily beneficial.

\question{R4.Q4}{Can Lean's motive computation be called directly with useful failures?}
\status{Yes. Lean 4.34.0 exposes concrete meta-level entry points; their cost and diagnostic abstraction must be controlled by an adapter.}

Three operations in the inspected source are directly relevant. The function
\code{getMajorTypeIndices} extracts the recursor's indices. The function
\code{mkRecursorAppPrefix} constructs a recursor prefix through its motive by
abstracting the major premise and indices from the target. Finally, the goal
method \code{MVarId.induction} already reverts indices and the major premise,
reintroduces them, builds the prefix, and returns branch goals with a
free-variable substitution. These operations are not all hidden inside syntax
elaboration \cite{leaninduction}.

The low-level index extractor requires variable indices satisfying specific
occurrence and dependency conditions. A major premise with a compound or
repeated index may therefore need the higher-level preparation used by
\code{cases}/\code{induction}, or an explicit generalization with equality
constraints. Calling the low-level prefix builder without preparing its
preconditions is not a complete motive-inference algorithm.

Implement a version-pinned adapter returning a structured result:
\begin{lstlisting}
MotiveAttempt :=
  success(prefix, branches, localSubstitution, dependencies)
  | unsupportedMajor(reason)
  | indexPreparationNeeded(indices)
  | typeMismatch(diagnostic)
  | unresolvedConstraint(constraint)
  | resourceExhausted
\end{lstlisting}
Run every attempt in a saved state with explicit heartbeat and recursion
limits. Restore assignments, messages, local substitutions, and branch-local
caches on failure. Do not parse fragile exception text as a proof of
impossibility. Preserve the original diagnostic for developers while exposing
a stable coarse category to the search.

This makes the initial indexed-induction improvement relatively contained:
remove Leant2's blanket exclusion for a well-prepared subset of indexed
families, reuse Lean's construction, and test the resulting branches. It does
not solve general higher-order motive search, but it is a useful prerequisite
before investing in a new motive language.
